\scanpage{119}
% archive sheet number: 118

2) $\otimes$-Foncteur $F:C\handwritten{\longrightarrow}C'$ (avec $C,C'$ des trucs-catégories).

Les données sont un foncteur \deleted{add} $F$, plus \deleted{\unsure{deu|deux}}un isomorphisme\deleted{s} de \unsure{bifoncteurs|bifoncteur}
\[
c(X,Y):F(X\otimes Y)\handwritten{\xrightarrow{\sim}}F(X)\otimes F(Y)
\]
\deleted{et isom.}
\[
\deleted{i:F(\underline{1}_C)\unsure{\longrightarrow|\simeq}\unsure{\underline{1}_{C'}|\underline{1}'_C}}
\]
satisfaisant les conditions:

a) Compatibilité \inserted{de $c$} avec la donnée de $\handwritten{\phi,\phi'}$

b) Compatibilité \inserted{de $c$} avec la donnée de $\handwritten{\unsure{\mu|\psi},\mu'}$

c) \deleted{Compatibilité de $i$ \unsure{avec} et $c$}\inserted{Existe isom.\ \deleted{fonct} $i\handwritten{\unsure{:|(}}F(\underline{1})\handwritten{\simeq}\underline{1}'$ compatible avec $c$ et} avec les isomorphismes $u\inserted{u'}$ (ou ce \deleted{\unsure{qui}} qui revient au même, avec \inserted{$c$ et} $v,v'$).

NB Cet $i$ sera alors uniquement déterminé.

d) $F$ est additif.

3) Homomorphisme d'un $\otimes$-foncteur $F$ dans un autre $G$: c'est un homomorphisme de foncteurs $f:F\handwritten{\longrightarrow}G$ satisfaisant le\deleted{s} condition

a) compatibilité avec \deleted{$c$ et} $c_F,c_G$

\handwritten{\deleted{\unsure{(A)|(a)} \unsure{pour} \illegible{2}}}

\deleted{b)} \handwritten{\deleted{\unsure{(A)|(a)}}(}compatibilité avec $i_F,i_G$ (\deleted{c'}\handwritten{\unsure{ce n'|car}} est \deleted{peut-être}\inserted{\handwritten{\deleted{\unsure{sans doute}} sans doute \unsure{pas une|par un}}} conséquence formelle de a) - à regarder\deleted{)}\handwritten{; cf plus bas)}

\handwritten{4)} \uline{Remarque\handwritten{s}}\qquad A dissocier \inserted{\handwritten{ce qui précède}} en plusieurs aspects de la donnée d'un bifoncteur\handwritten{$\otimes$}:

a) Donnée d'associativité, soumise à l'axiome du pentagone. Un foncteur ``compatible avec $\otimes$\inserted{associat\deleted{t}if''} de $C$ dans $C'$ est un foncteur $+$ un $c$ , compatible avec la donnée de $\handwritten{\phi}$ .

b) Donnée d'unité bilatère $\underline{1},u,v$ (unique à isomorphisme unique près). Un foncteur ``compatible avec \inserted{\handwritten{$\otimes$ à}} unité\deleted{s} bilatère\deleted{s}'' est un $F$ muni de $c$, tel qu'il existe $i$ compatible avec $c,u,v$ et $c',u',v'$ (lequel $i$ sera alors unique).

\margin{\handwritten{\uline{NB} Le monoïde des endom. de $\underline{1}$ est alors comm.}}

\handwritten{c)} Donnée de commutativité $\handwritten{\mu}$ , satisfaisant la condition de réflexivité.\deleted{(} Un foncteur
